\documentclass[11pt]{article}

\usepackage[margin=1in]{geometry}
\usepackage[T1]{fontenc}
\usepackage[utf8]{inputenc}
\usepackage{lmodern}
\usepackage{microtype}
\usepackage{hyperref}
\usepackage{booktabs}
\usepackage{longtable}
\usepackage{tabularx}
\usepackage{array}
\usepackage{enumitem}
\usepackage[table]{xcolor}
\usepackage{titlesec}
\usepackage{fancyhdr}
\usepackage{lastpage}
\usepackage[most]{tcolorbox}

\definecolor{EvalOpsBlue}{HTML}{0F4C81}
\definecolor{EvalOpsInk}{HTML}{13212B}
\definecolor{EvalOpsSlate}{HTML}{5A6B7A}
\definecolor{EvalOpsMist}{HTML}{EEF4F8}
\definecolor{EvalOpsLine}{HTML}{D5E0E8}
\definecolor{EvalOpsAccent}{HTML}{2C7DA0}

\hypersetup{
  colorlinks=true,
  linkcolor=EvalOpsBlue,
  urlcolor=EvalOpsAccent,
  pdftitle={EvalOps Research Report: Cognitive Dissonance DSPy},
  pdfauthor={Jonathan Haas},
  pdfsubject={Empirical study of proof-oriented conflict resolution in multi-agent LLM systems},
  pdfcreator={Tectonic}
}

\setlist[itemize]{leftmargin=1.4em, itemsep=0.25em, topsep=0.4em}
\setlist[enumerate]{leftmargin=1.7em, itemsep=0.25em, topsep=0.4em}
\setlength{\parskip}{0.55em}
\setlength{\parindent}{0pt}
\setlength{\emergencystretch}{2em}
\renewcommand{\arraystretch}{1.15}

\titleformat{\section}
  {\Large\bfseries\color{EvalOpsBlue}}
  {\thesection}
  {0.75em}
  {}

\titleformat{\subsection}
  {\large\bfseries\color{EvalOpsInk}}
  {\thesubsection}
  {0.65em}
  {}

\pagestyle{fancy}
\fancyhf{}
\fancyhead[L]{\small\color{EvalOpsSlate}\textbf{EvalOps Research}}
\fancyhead[R]{\small\color{EvalOpsSlate}Cognitive Dissonance DSPy}
\fancyfoot[C]{\small\color{EvalOpsSlate}\thepage\ of \pageref{LastPage}}
\renewcommand{\headrulewidth}{0.4pt}
\renewcommand{\headrule}{\hbox to\headwidth{\color{EvalOpsLine}\leaders\hrule height \headrulewidth\hfill}}
\renewcommand{\footrulewidth}{0pt}

\newtcolorbox{summarybox}{
  colback=EvalOpsMist,
  colframe=EvalOpsBlue,
  boxrule=0.8pt,
  arc=2pt,
  left=10pt,
  right=10pt,
  top=8pt,
  bottom=8pt
}

\newtcolorbox{calloutbox}{
  colback=white,
  colframe=EvalOpsLine,
  boxrule=0.7pt,
  arc=2pt,
  left=10pt,
  right=10pt,
  top=8pt,
  bottom=8pt
}

\begin{document}
\hypersetup{pageanchor=false}
\pagenumbering{roman}

\begin{titlepage}
\thispagestyle{empty}

\vspace*{1.5cm}

\begin{center}
{\Large\color{EvalOpsSlate}\textbf{EVALOPS RESEARCH TECHNICAL REPORT}\par}
\vspace{0.65cm}
{\Huge\bfseries\color{EvalOpsInk}Cognitive Dissonance DSPy\par}
\vspace{0.2cm}
{\Large\color{EvalOpsBlue}An Empirical Study of Proof-Oriented Conflict Resolution in Multi-Agent LLM Systems\par}
\end{center}

\vspace{0.8cm}

\color{EvalOpsLine}\rule{\linewidth}{1.2pt}

\vspace{0.7cm}

\begin{tabularx}{\linewidth}{@{}lX@{}}
\textbf{Author} & Jonathan Haas \\
\textbf{Email} & \texttt{jonathan@evalops.dev} \\
\textbf{Organization} & EvalOps \\
\textbf{Prepared} & 2026-04-07 \\
\textbf{Study source revision} & \texttt{d11d492-dirty} \\
\textbf{Document type} & Experimental technical report \\
\end{tabularx}

\vspace{0.8cm}

\begin{summarybox}
\textbf{Study objective.} This report evaluates whether the repository can move
agent disagreement from debate into proof. The work reported here is based on
actual experiments run from \nolinkurl{research/run_study.py}, not on static
repository claims.
\end{summarybox}

\vfill

\begin{calloutbox}
\textbf{Suggested citation}\\
Haas, Jonathan. \emph{Cognitive Dissonance DSPy: An Empirical Study of
Proof-Oriented Conflict Resolution in Multi-Agent LLM Systems}. EvalOps
Research Technical Report, 2026-04-07.
\end{calloutbox}

\vfill

{\small\color{EvalOpsSlate}
This report is grounded in repository-local experiments, generated study
artifacts, local test runs, and a live OpenAI-compatible extraction benchmark
executed on 2026-04-07.}

\end{titlepage}

\begin{abstract}
This report evaluates the current state of Cognitive Dissonance DSPy as a
research artifact. The central question is whether a multi-agent LLM system can
resolve formalizable disagreements through proof instead of through debate or
confidence aggregation. To answer that question, I ran a curated symbolic
verification benchmark, a live natural-language extraction benchmark against an
OpenAI-compatible endpoint, and a small stability check.

The main result is that the current pipeline saturates the repository's
internal benchmark once Coq is available locally. Direct pattern matching
translates only 25\% of the formalizable natural-language cases, while the
current extraction pipeline reaches 100\% canonicalization accuracy on the same
set. With \texttt{coqc} and \texttt{coqchk} installed, the formal benchmark
becomes 16/16 decisive and the formalizable slice of the live extraction
benchmark becomes 12/12 decisive. Six of those 12 extraction cases are
machine-checked locally, while the remainder resolve through SMT proofs or
decisive refutations. The strongest conclusion is therefore no longer ``proof
is still the bottleneck on this benchmark,'' but rather ``the benchmark is no
longer hard enough to be the main source of uncertainty.''
\end{abstract}

\section*{Executive Summary}

\begin{summarybox}
\begin{itemize}
  \item \textbf{Direct finding.} On the curated extraction benchmark, direct
  translation succeeds on only 3 of 12 formalizable natural-language claims,
  while the extraction pipeline produces exact canonical forms for all 12.
  \item \textbf{Coq-enabled result.} With local \texttt{coqc} and
  \texttt{coqchk} available, the symbolic benchmark reaches 16/16 decisive
  outcomes and the formalizable extraction slice reaches 12/12 decisive
  outcomes.
  \item \textbf{Evidence mix.} Those decisive outcomes are not all equivalent:
  the formal benchmark contains 3 machine-checked cases, while the extraction
  benchmark contains 6 machine-checked cases, 3 SMT proofs, and 3 decisive
  refutations.
  \item \textbf{Research implication.} The system's key value today is honest
  evidence typing: SMT-backed, derived, machine-checked, and inconclusive
  outcomes are kept distinct rather than merged into a single ``proved'' label.
\end{itemize}
\end{summarybox}

\tableofcontents
\newpage
\hypersetup{pageanchor=true}
\pagenumbering{arabic}

\section{Objective}

Most multi-agent systems handle disagreement with prompts, debate, voting, or
confidence aggregation. Those strategies are appropriate when the disagreement
is subjective. They are weaker when the disputed claim is formalizable.

This repository explores a different policy:

\begin{quote}
If a claim can be translated into a formal object, attempt proof before
arbitration.
\end{quote}

The goal of this study was not to restate the architecture, but to run real
experiments and determine what the system currently does, what it does not yet
do, and where the next research effort should go.

\section{Research Questions}

The study focused on three practical questions.

\begin{enumerate}
  \item \textbf{RQ1:} Does necessity-enhanced verification materially improve
  symbolic claim handling?
  \item \textbf{RQ2:} Does the OpenAI-compatible extraction pipeline improve
  translation success over direct pattern matching on natural-language claims?
  \item \textbf{RQ3:} After successful extraction, how much of the remaining
  problem is still in proof rather than extraction?
\end{enumerate}

\section{Experimental Setup}

\subsection{Code and Provenance}

All study artifacts were generated by \nolinkurl{research/run_study.py}. The
study metadata recorded the source revision as \texttt{d11d492-dirty}, meaning
the experiments were run on top of commit \texttt{d11d492} with local
uncommitted modifications that are part of the final repository state produced
by this work. The exact experiment outputs are stored in:

\begin{itemize}
  \item \nolinkurl{research/results/study_results.json}
  \item \nolinkurl{research/results/study_summary.md}
\end{itemize}

\subsection{Local Environment}

\begin{longtable}{>{\raggedright\arraybackslash}p{0.32\linewidth} >{\raggedright\arraybackslash}p{0.58\linewidth}}
\toprule
\textbf{Item} & \textbf{Observed configuration} \\
\midrule
Python & \texttt{3.12.8} \\
SMT solver & \texttt{z3-solver} installed locally \\
Coq local availability & \texttt{coqc} and \texttt{coqchk} available \\
Live extraction provider & OpenRouter-compatible endpoint \\
Model & \texttt{openai/gpt-4.1-mini} \\
Extraction temperature & \texttt{0.0} \\
Study date & 2026-04-07 \\
\bottomrule
\end{longtable}

This matters because the local study now exercises the full proof pipeline
rather than falling back to an SMT-only or no-Coq configuration.

\subsection{Evidence Policy}

The study intentionally separates \emph{decisive} and \emph{optimistic}
outcomes.

\begin{itemize}
  \item \textbf{Decisive:} statuses strong enough to resolve truth or falsity
  directly, such as \texttt{smt\_proved} and \texttt{smt\_refuted}.
  \item \textbf{Optimistic:} statuses that include derived support such as
  \texttt{derived\_proved} and \texttt{derived\_refuted}. These do not count as
  decisive evidence.
\end{itemize}

This distinction is not bookkeeping. It is the central epistemic discipline of
the project.

\section{Benchmarks}

\subsection{Formal Verification Benchmark}

The symbolic benchmark contains 16 claims across seven families:

\begin{itemize}
  \item arithmetic
  \item multiplication
  \item subtraction
  \item factorial
  \item Fibonacci
  \item GCD
  \item inequalities
\end{itemize}

Each family includes true and false examples. The benchmark compares two
conditions:

\begin{itemize}
  \item hybrid verification with necessity enabled
  \item hybrid verification without necessity
\end{itemize}

\subsection{Extraction Benchmark}

The natural-language benchmark contains 16 cases:

\begin{itemize}
  \item 12 formalizable utterances
  \item 4 unformalizable subjective controls
\end{itemize}

The extraction study compares:

\begin{itemize}
  \item direct translation from raw natural language
  \item the structured extraction pipeline followed by translation and proof
\end{itemize}

\subsection{Stability Check}

A small repeated-measures check runs extraction twice on four representative
cases:

\begin{itemize}
  \item arithmetic true
  \item factorial true
  \item implication
  \item subjective control
\end{itemize}

This is not a full calibration study. It is only a basic check for
deterministic behavior at temperature 0.0.

\section{Results}

\subsection{RQ1: Necessity and Symbolic Verification}

\begin{center}
\begin{tabular}{lrrrr}
\toprule
\textbf{Condition} & \textbf{Cases} & \textbf{Decisive cov.} & \textbf{Decisive acc.} & \textbf{Optimistic acc.} \\
\midrule
Hybrid + necessity & 16 & 100.0\% & 100.0\% & 100.0\% \\
Hybrid without necessity & 16 & 100.0\% & 100.0\% & 100.0\% \\
\bottomrule
\end{tabular}
\end{center}

Mean proof time was approximately 154.0\,ms in both conditions on this
benchmark. Both conditions also produced 3 machine-checked positive cases:
factorial, Fibonacci, and GCD.

\textbf{Interpretation.} With Coq available locally, the symbolic benchmark is
fully decisive in both conditions. Necessity therefore has no measurable effect
on headline accuracy or coverage on this suite. That is not evidence that
necessity is unhelpful in general; it is evidence that this benchmark no longer
separates the two conditions.

\subsection{RQ2: Extraction Versus Direct Translation}

\begin{center}
\begin{tabular}{lr}
\toprule
\textbf{Metric} & \textbf{Result} \\
\midrule
Formalizable natural-language cases & 12 \\
Direct translator success & 25.0\% \\
Extractor formalizability accuracy (all 16 cases) & 100.0\% \\
Extractor exact canonical match (formalizable cases) & 100.0\% \\
Translation success after extraction & 100.0\% \\
End-to-end decisive coverage after extraction & 100.0\% \\
Machine-checked formalizable cases after extraction & 6 \\
Mean extraction latency & 2747\,ms/case \\
\bottomrule
\end{tabular}
\end{center}

\textbf{Interpretation.} On the curated benchmark, the extraction pipeline now
completely solves the canonicalization problem. Every formalizable
natural-language case was converted into the exact canonical claim expected by
the benchmark, all four subjective controls were rejected, and every
formalizable case reached a decisive downstream result.

\subsection{RQ3: Proof After Successful Extraction}

After perfect canonicalization on the benchmark, decisive proof coverage no
longer remains limited:

\begin{itemize}
  \item end-to-end decisive coverage on formalizable extraction cases:
  100.0\%
  \item end-to-end decisive accuracy on formalizable extraction cases:
  100.0\%
\end{itemize}

In concrete terms, all 12 formalizable natural-language cases were resolved
decisively after extraction. The evidence mix is:

\begin{itemize}
  \item \textbf{machine-checked:} factorial, Fibonacci, GCD, implication,
  universal quantification, and existential quantification
  \item \textbf{SMT-proved:} arithmetic true, multiplication true, and
  inequality true
  \item \textbf{decisive refutations:} arithmetic false, subtraction false,
  and inequality false
\end{itemize}

\textbf{Interpretation.} On the current internal suite, proof is no longer the
measured bottleneck once Coq is installed locally. The remaining limitation is
that this benchmark is now too easy relative to the repository's current
capability.

\subsection{Stability}

The repeated extraction check reached a 100.0\% stable-case rate over 4 cases
and 2 trials per case. This is a small result, but it suggests that the
pipeline is operationally stable under deterministic decoding on the chosen
provider and model.

\section{Findings}

\subsection{Finding 1: Extraction Was the First Bottleneck}

The benchmark data shows a clean before-and-after story. Direct translation from
raw natural language succeeded on only 3 of 12 formalizable cases. The
extraction pipeline pushed that to 12 of 12. That closes one important gap.

With local Coq available, however, the end-to-end problem disappears on this
suite. That means the benchmark has become a confirmation harness rather than a
stress test.

\subsection{Finding 2: This Benchmark No Longer Separates Necessity from Non-Necessity}

Both symbolic benchmark conditions now reach 100\% decisive coverage and 100\%
decisive accuracy. That means this benchmark can no longer tell us whether
necessity provides unique value under harder conditions. A stronger ablation
needs harder recursive and logical cases.

\subsection{Finding 3: Evidence Strength Still Varies Even When Accuracy Saturates}

The current extraction benchmark resolves all 12 formalizable cases decisively,
but only 6 of those are machine-checked. The rest are SMT proofs or decisive
refutations. That distinction matters because a 100\% headline number can hide
different levels of evidential strength.

\subsection{Finding 4: The Proof-Status Model Is Working as Intended}

The repository does not flatten derived support into decisive proof. That
decision now has experimental value: without it, the study would overstate the
system's capability. With it, the study can say something precise:

\begin{quote}
This benchmark is solved, but only part of it is machine-checked.
\end{quote}

\section{Threats to Validity}

\subsection{Small Internal Benchmark}

The benchmark set is small and authored within the project. That makes it
useful for engineering truth, but weak for strong claims about generalization.

\subsection{Local Coq Availability Changes the Result Materially}

Installing Coq locally turned the earlier partial results into full decisive
coverage on this benchmark. That is useful, but it also means any report that
ignores environment configuration would be misleading. The benchmark should be
interpreted together with the local proof-toolchain state.

\subsection{Provider Dependence}

The live extraction results depend on one provider configuration on one date:
OpenRouter-compatible routing with \texttt{openai/gpt-4.1-mini} at temperature
0.0. Different models or provider settings could behave differently.

\subsection{Benchmark-Aware Improvement}

This work used the observed failures to improve the system and rerun the study.
That is the correct engineering loop, but it also means the final benchmark
results are partially the result of targeted iteration against the benchmark's
failure modes.

\section{Reproducibility}

The study can be reproduced from repository-local artifacts.

Current local validation status after installing Coq:
\texttt{248 passed, 1 skipped} from
\texttt{.venv/bin/python -m pytest -q}, plus a passing live provider-backed
integration test against the configured OpenAI-compatible endpoint.

\subsection{Commands}

\begin{calloutbox}
\texttt{.venv/bin/python research/run\_study.py}
\end{calloutbox}

\begin{calloutbox}
\texttt{OPENAI\_API\_KEY=... OPENAI\_BASE\_URL=... OPENAI\_MODEL=...}\\
\texttt{.venv/bin/python research/run\_study.py --temperature 0.0 --stability-trials 2}
\end{calloutbox}

\begin{calloutbox}
\texttt{.venv/bin/python -m pytest -q}
\end{calloutbox}

\subsection{Artifacts}

\begin{itemize}
  \item \nolinkurl{research/benchmarks/formal_verification_benchmark.json}
  \item \nolinkurl{research/benchmarks/extraction_benchmark.json}
  \item \nolinkurl{research/results/study_results.json}
  \item \nolinkurl{research/results/study_summary.md}
\end{itemize}

\section{Conclusion}

This repository now supports a defensible experimental claim:

\begin{quote}
On a small internal benchmark with Coq installed locally, Cognitive Dissonance
DSPy can canonicalize formalizable natural-language disagreements and resolve
them decisively end to end.
\end{quote}

That is stronger than the earlier no-Coq result, but it does not mean the
general problem is solved. The next research task is no longer ``fix the
current benchmark,'' but ``build a benchmark that is difficult enough to
measure where the pipeline still fails, while continuing to report evidence
strength honestly.''

\section{References}

\begin{thebibliography}{99}
\small

\bibitem{study_runner}
EvalOps. \emph{Study runner}. Repository artifact,
\nolinkurl{research/run_study.py}, source revision \texttt{d11d492-dirty},
2026.

\bibitem{study_results}
EvalOps. \emph{Generated study results}. Repository artifact,
\nolinkurl{research/results/study_results.json}, source revision
\texttt{d11d492-dirty}, 2026.

\bibitem{study_summary}
EvalOps. \emph{Generated study summary}. Repository artifact,
\nolinkurl{research/results/study_summary.md}, source revision
\texttt{d11d492-dirty}, 2026.

\bibitem{readme}
EvalOps. \emph{Repository README}. Repository artifact, \nolinkurl{README.md},
source revision \texttt{d11d492-dirty}, 2026.

\bibitem{types}
EvalOps. \emph{Proof status model}. Repository artifact,
\nolinkurl{formal_verification/types.py}, source revision
\texttt{d11d492-dirty}, 2026.

\bibitem{detector}
EvalOps. \emph{Formal verification detector}. Repository artifact,
\nolinkurl{formal_verification/detector.py}, source revision
\texttt{d11d492-dirty}, 2026.

\bibitem{openai_agents}
EvalOps. \emph{OpenAI-compatible extraction pipeline}. Repository artifact,
\nolinkurl{formal_verification/openai_agents.py}, source revision
\texttt{d11d492-dirty}, 2026.

\bibitem{necessity_prover}
EvalOps. \emph{Necessity-based proof layer}. Repository artifact,
\nolinkurl{formal_verification/necessity_prover.py}, source revision
\texttt{d11d492-dirty}, 2026.

\bibitem{translator}
EvalOps. \emph{Claim translator}. Repository artifact,
\nolinkurl{formal_verification/translator.py}, source revision
\texttt{d11d492-dirty}, 2026.

\bibitem{z3_prover}
EvalOps. \emph{Hybrid and SMT proving}. Repository artifact,
\nolinkurl{formal_verification/z3_prover.py}, source revision
\texttt{d11d492-dirty}, 2026.

\end{thebibliography}

\end{document}
